\documentclass[11pt,a4paper]{article}
\usepackage[spanish,es-tabla]{babel}
\usepackage[utf8]{inputenc}
\usepackage[T1]{fontenc}
\usepackage{lmodern}
\usepackage[margin=2.2cm]{geometry}
\usepackage{graphicx,booktabs,multirow,float,amsmath,amssymb,listings,xcolor,hyperref}
\usepackage{caption}
\usepackage{pgfplots}
\pgfplotsset{compat=1.17, /pgf/number format/use period}
\captionsetup{font=small,labelfont=bf}
\setlength{\emergencystretch}{2em}
\graphicspath{{figs/}}
\lstset{basicstyle=\ttfamily\footnotesize,frame=single,breaklines=true,
        keywordstyle=\color{blue!60!black},commentstyle=\color{gray}}

% Notas del borrador: borrar \nota{...} al redactar la version final.
\newcommand{\nota}[1]{\par\noindent\colorbox{yellow!25}{\parbox{\dimexpr\linewidth-2\fboxsep}{\small\textbf{Borrador:} #1}}\par}
% Datos que faltan: reemplazar \falta{...} por el valor medido.
\newcommand{\falta}[1]{\textcolor{red}{\textbf{[#1]}}}
\newcommand{\Tcomp}{T_{\text{comp}}}
\newcommand{\Tcomm}{T_{\text{comm}}}
\newcommand{\Ttot}{T_{\text{total}}}

\title{Tarea 3: Factorización LU distribuida con MPI4PY\\
\large IIC3533 Computación de Alto Rendimiento $\cdot$ 2026-2}
\author{Integrante 1 \and Integrante 2 \and Integrante 3}
\date{16 de noviembre de 2026}

\begin{document}
\decimalpoint
\maketitle

% ======================================================================
\section{Introducción}
La factorización $A=LU$ de una matriz densa $A\in\mathbb{R}^{N\times N}$ cuesta $\tfrac{2}{3}N^3$ FLOP y es la base de los métodos directos para resolver sistemas lineales. En esta tarea se implementa la versión por bloques de $b\times b$ sobre un \emph{grid} 2D de $p_r\times p_c$ procesos MPI con distribución bloque-cíclica, usando \texttt{mpi4py} para la comunicación y NumPy/SciPy (BLAS/LAPACK) para el cómputo local. Se estudia el efecto del tamaño de bloque $b$, del número de procesos y de \emph{threads} por proceso, de la forma del \emph{grid} y de la ubicación de los procesos (uno o dos nodos), separando en cada caso el tiempo de cómputo y el de comunicación. Al final se contrastan las mediciones con el modelo de costo $(\alpha,\beta,\gamma)$.

Como la matriz es diagonal dominante, la factorización se completa sin pivoteo, por lo que todo el trabajo se reduce a tres operaciones por paso $k$: factorizar $A_{kk}$ (\texttt{lu\_factor}), calcular los paneles $U_{kj}=L_{kk}^{-1}A_{kj}$ y $L_{ik}=A_{ik}U_{kk}^{-1}$ (\texttt{solve\_triangular}) y actualizar $A_{ij}\leftarrow A_{ij}-L_{ik}U_{kj}$ (\texttt{@}).

% ======================================================================
\section{Metodología}
\subsection{Plataformas}
\begin{table}[H]
\centering\small
\caption{Características del computador personal (parte a y b) y del \emph{cluster} Kraken (partes c--f).}
\label{tab:maquinas}
\begin{tabular}{lp{5.2cm}p{6.2cm}}
\toprule
 & \textbf{Computador personal} & \textbf{Cluster Kraken} \\
\midrule
CPU & Apple M5 (ARM64) & \falta{modelo CPU (\texttt{lscpu})} \\
Cores & 10: 4 de rendimiento (P) + 6 de eficiencia (E), sin SMT & 40 CPUs por nodo; 2 nodos (\texttt{ventress}, \texttt{yodaxico}) \\
Cachés & L1d 128 KB (P) / 64 KB (E); L2 16 MB (clúster P) + 6 MB (clúster E) & \falta{cachés} \\
Memoria & 16 GB unificada & $\sim$192 GB por nodo \\
Límites del grupo & -- & 8 CPUs y 4 GB por nodo, 15 min por \emph{job} \\
Red entre nodos & -- & \falta{Ethernet / InfiniBand} \\
Sistema operativo & macOS 26.6 & \falta{SO del cluster} \\
Python / MPI & Python 3.12, NumPy 2.5, SciPy 1.18, mpi4py 4.1, Open MPI 5.0 & Miniconda (\texttt{hpc\_mpi}), \falta{Open MPI x.y, mpi4py x.y} \\
BLAS & Apple Accelerate & \falta{OpenBLAS / MKL} \\
\bottomrule
\end{tabular}
\end{table}
\nota{Completar la columna del \emph{cluster} con la salida de \texttt{info\_nodo} (al inicio de \texttt{resultados/t3\_c\_f\_*.out} y \texttt{t3\_e\_*.out}): modelo de CPU, cachés, SO, versiones, BLAS y tipo de red (importante para la parte e).}

\subsection{Implementación y medición}
\begin{itemize}
  \item \textbf{Matriz de prueba.} El bloque $(I,J)$ se genera con un generador cuya semilla depende solo de $(I,J)$, con entradas uniformes en $(-1,1)$ y diagonal igual a $N$. Cada fila tiene $N-1$ entradas fuera de la diagonal con $|a_{ij}|<1$, así que $\sum_{j\ne i}|a_{ij}|<N=a_{ii}$: la matriz es estrictamente diagonal dominante y la LU existe sin pivoteo. Como cada bloque depende solo de su posición, cada proceso genera sus propios bloques sin comunicación (no hay que armar la matriz completa en el rank 0, lo que con $N=8192$ ocuparía 512 MB adicionales), y la versión serial genera exactamente la misma matriz. La generación no se incluye en los tiempos.
  \item Los tiempos se miden con \texttt{MPI.Wtime()} después de un \texttt{Barrier}. $\Tcomp$ acumula \texttt{lu\_factor}, \texttt{solve\_triangular} y la actualización; $\Tcomm$ acumula los \texttt{Bcast} y el empaquetado de los paneles en buffers contiguos; $\Ttot$ va hasta el \texttt{Barrier} final. Como \texttt{Bcast} bloquea hasta que llega el mensaje, $\Tcomm$ incluye también el tiempo que un proceso \emph{espera} a la raíz (desbalance de carga), no solo la transferencia. Cada proceso registra sus propios tiempos y el rank 0 los reúne con \texttt{gather}; para cada corrida se reporta el máximo sobre los procesos.
  \item Se reporta el rendimiento como GFLOP/s $=\tfrac{2}{3}N^3/\Ttot$.
  \item En el \emph{cluster}, \texttt{OMP\_NUM\_THREADS}, \texttt{OPENBLAS\_NUM\_THREADS} y \texttt{MKL\_NUM\_THREADS} se fijan en $T$ en cada corrida y se propagan a todos los procesos (\texttt{mpirun -x}). El número de procesos por nodo se controla con \texttt{-{}-ntasks-per-node} y \texttt{-{}-map-by ppr:$k$:node}, y en cada corrida se usan exactamente $p\cdot T\le8$ CPUs por nodo, sin sobresuscripción. Se usa \texttt{-{}-bind-to none} para que los $T$ threads de un proceso no queden fijos en un solo core.
  \item Cada configuración se ejecuta 3 veces dentro del mismo \emph{job} y se reporta la \textbf{mediana}, como en la Tarea 2.
\end{itemize}

% ======================================================================
\section{(a) LU serial y tamaño de bloque}
Se ejecutó \texttt{lu\_serial.py} con $N=2^{12}=4096$ en el computador personal con un \emph{thread} (\texttt{OMP\_NUM\_THREADS}, \texttt{OPENBLAS\_NUM\_THREADS}, \texttt{MKL\_NUM\_THREADS} y \texttt{VECLIB\_MAXIMUM\_THREADS} en 1), variando $b$ de $2^{12}$ a $2^6$, con 3 repeticiones por valor (\texttt{run\_experimento\_a.py}).

\begin{figure}[H]\centering
\begin{tikzpicture}
\begin{axis}[
  width=0.75\linewidth, height=6cm,
  xmode=log, log basis x=2,
  xlabel={Tamaño de bloque $b$}, ylabel={Tiempo [s]},
  title={LU serial, $N=4096$, 1 \emph{thread} (mediana de 3)},
  xtick={64,128,256,512,1024,2048,4096},
  xticklabels={64,128,256,512,1024,2048,4096},
  ymin=0, ymax=0.4, grid=major, label style={font=\small}, tick label style={font=\small},
  yticklabel style={/pgf/number format/fixed, /pgf/number format/precision=2},
  nodes near coords, point meta=explicit symbolic,
  every node near coord/.append style={font=\scriptsize, anchor=south, yshift=2pt, fill=white, inner sep=1pt}]
\addplot[mark=*, thick, blue!70!black] coordinates {
  (4096,0.2684)[0.268] (2048,0.2205)[0.221] (1024,0.1947)[0.195]
  (512,0.1747)[0.175] (256,0.2082)[0.208] (128,0.2567)[0.257] (64,0.3316)[0.332]};
\end{axis}
\end{tikzpicture}
\caption{Tiempo de \texttt{lu\_serial.py} con $N=4096$ según $b$. $b=4096$ corresponde a una sola llamada a \texttt{lu\_factor} (LAPACK).}
\label{fig:a}
\end{figure}

\begin{table}[H]
\centering\small
\caption{LU serial con $N=4096$ y un \emph{thread}. Tiempo: mediana y rango de 3 repeticiones. La columna ``\% en \texttt{@}'' es la fracción de los FLOP que corresponde a la actualización (producto de matrices); el resto es \texttt{lu\_factor} y \texttt{solve\_triangular}.}
\label{tab:a}
\begin{tabular}{rrrrrr}
\toprule
$b$ & $N/b$ & \% en \texttt{@} & Tiempo [s] & Rango [s] & GFLOP/s \\
\midrule
4096 & 1  & 0    & 0.268 & 0.267--0.270 & 171 \\
2048 & 2  & 37.5 & 0.221 & 0.220--0.228 & 208 \\
1024 & 4  & 65.6 & 0.195 & 0.186--0.197 & 235 \\
\textbf{512} & \textbf{8} & \textbf{82.0} & \textbf{0.175} & \textbf{0.174--0.186} & \textbf{262} \\
256  & 16 & 90.8 & 0.208 & 0.192--0.231 & 220 \\
128  & 32 & 95.4 & 0.257 & 0.255--0.264 & 178 \\
64   & 64 & 97.7 & 0.332 & 0.331--0.348 & 138 \\
\bottomrule
\end{tabular}
\end{table}

El mejor tiempo se obtiene con $\boldsymbol{b=512}$ (0.175 s, 262 GFLOP/s); $b=1024$ queda a un 11\% y $b=256$ a un 19\%. La curva tiene forma de U, con un factor 1.9 entre el mejor ($b=512$) y el peor caso ($b=64$).

\textbf{Rendimiento absoluto.} 262 GFLOP/s con un solo \emph{thread} es muy alto para un core. La razón es que Accelerate, la BLAS/LAPACK de Apple que usa NumPy en macOS, ejecuta los productos de matrices en la unidad de matrices del procesador. Se comprobó que efectivamente corre en un solo \emph{thread}: con $N=8192$, el tiempo de CPU (3.05 s) es prácticamente igual al tiempo real (3.26 s), mientras que sin limitar los \emph{threads} el tiempo de CPU supera al real (4.00 frente a 2.80 s). Por esto los valores absolutos no son comparables con los del \emph{cluster}, pero la forma de la curva sí se explica con los mismos argumentos.

\textbf{$b$ grande ($\ge1024$).} Con pocos bloques, una parte importante del trabajo no es producto de matrices. En total, la factorización de los bloques diagonales suma $\tfrac{2}{3}Nb^2$ FLOP y los paneles $\sim N^2b$, frente a $\tfrac{2}{3}N^3$ del total, y la Tabla~\ref{tab:a} muestra que con $b=2048$ solo el 37.5\% de los FLOP están en el \texttt{@}. \texttt{lu\_factor} (\texttt{getrf}) y \texttt{solve\_triangular} (\texttt{trsm}) tienen dependencias internas y rinden menos que el producto de matrices: el caso extremo $b=4096$ es una sola llamada a LAPACK y alcanza 171 GFLOP/s, frente a 262 con $b=512$.

\textbf{$b$ pequeño ($\le256$).} Casi todo el trabajo es la actualización, pero es un producto de rango $b$: en cada paso se lee y escribe toda la submatriz restante (hasta 128 MB, mucho más que los 16 MB de L2) para hacer solo $2b$ FLOP por elemento. La intensidad aritmética es $\sim b/8$ FLOP por byte y el tráfico total con memoria crece como $\tfrac{16N^3}{3b}$ bytes: $\approx5.7$ GB con $b=64$ frente a $\approx0.7$ GB con $b=512$. Además, la unidad de matrices es menos eficiente con una dimensión interna pequeña, y hay $N/b$ pasos con llamadas pequeñas a LAPACK y el costo fijo de Python.

\textbf{Óptimo.} $b=512$ equilibra ambos efectos: el 82\% de los FLOP están en el producto de matrices y la actualización tiene una intensidad de $\sim64$ FLOP/byte. Por lo tanto, se elige $\boldsymbol{b_a=512}$. Una primera medición del grupo en otro computador (Intel Pentium 4417U, con una versión anterior del código que actualizaba bloque a bloque) daba el óptimo en $b=128$: el $b$ óptimo depende del hardware (cachés, ancho de banda, BLAS) y de la implementación, lo que motiva repetir el barrido en el \emph{cluster} (parte c).

% ======================================================================
\section{(b) Implementación distribuida}
\subsection{Estructura}
\texttt{lu\_distributed.py} recibe \texttt{-{}-N}, \texttt{-{}-b}, \texttt{-{}-pr} y \texttt{-{}-pc}. Cada proceso calcula su posición en el \emph{grid} ($\textit{myrow}=\textit{rank}\,/\!/\,p_c$, $\textit{mycol}=\textit{rank}\bmod p_c$) y crea los comunicadores
\begin{lstlisting}[language=Python]
row_comm = comm.Split(myrow, mycol)
col_comm = comm.Split(mycol, myrow)
\end{lstlisting} Con estas claves, el rank de un proceso dentro de \texttt{row\_comm} es su columna y dentro de \texttt{col\_comm} es su fila, lo que simplifica la elección de la raíz. El proceso guarda los bloques $A_{IJ}$ con $I\equiv\textit{myrow}\pmod{p_r}$ y $J\equiv\textit{mycol}\pmod{p_c}$ en un arreglo local contiguo, ordenados por índice global.

\subsection{Comunicación en cada paso}
En el paso $k$, sea $(r,c)=(k\bmod p_r,\,k\bmod p_c)$ el dueño de $A_{kk}$. Todas las comunicaciones son \texttt{Bcast} (uno a todos), que es la colectiva natural porque un solo proceso tiene el dato y todos los de su fila o columna lo necesitan.

\begin{table}[H]
\centering\small
\caption{Operaciones colectivas en el paso $k$.}
\label{tab:colectivas}
\begin{tabular}{llll}
\toprule
\textbf{Dato} & \textbf{Colectiva} & \textbf{Comunicador} & \textbf{Raíz} \\
\midrule
$A_{kk}$ factorizado (hacia la fila $r$) & \texttt{Bcast} & \texttt{row\_comm} de la fila $r$ & $c$ \\
$A_{kk}$ factorizado (hacia la columna $c$) & \texttt{Bcast} & \texttt{col\_comm} de la columna $c$ & $r$ \\
Panel $L_{ik}$, $i>k$ & \texttt{Bcast} & \texttt{row\_comm} de cada fila & $c$ \\
Panel $U_{kj}$, $j>k$ & \texttt{Bcast} & \texttt{col\_comm} de cada columna & $r$ \\
\bottomrule
\end{tabular}
\end{table}

El bloque diagonal solo se envía a la fila $r$ (que lo necesita para calcular $U_{kj}=L_{kk}^{-1}A_{kj}$) y a la columna $c$ (que lo necesita para $L_{ik}=A_{ik}U_{kk}^{-1}$). Los paneles, en cambio, se difunden en \emph{todas} las filas y columnas del \emph{grid}, porque cada proceso necesita los $L_{ik}$ de sus filas de bloques y los $U_{kj}$ de sus columnas de bloques para actualizar su submatriz.

Como los bloques locales están ordenados por índice global, los bloques de un proceso con $I>k$ (o $J>k$) forman un sub-arreglo contiguo del arreglo local. Esto se aprovecha en tres lugares:
\begin{itemize}
  \item Cada proceso de la fila $r$ calcula \emph{todos} sus bloques $U_{kj}$ con una sola llamada a \texttt{solve\_triangular} sobre la franja $b\times(\text{columnas restantes})$, y los de la columna $c$ calculan todos sus $L_{ik}$ con otra llamada.
  \item Cada panel se difunde con \emph{un solo} \texttt{Bcast}: la raíz copia su franja a un buffer contiguo y el resto recibe en un buffer del mismo tamaño. Así hay 4 \texttt{Bcast} por paso como máximo, independientemente de cuántos bloques tenga cada proceso.
  \item La actualización es una sola multiplicación: \texttt{A[i0:, j0:] -= L\_panel @ U\_panel}.
\end{itemize}
Cada proceso genera sus propios bloques (ver Metodología), así que no hay distribución inicial de la matriz.

\subsection{Verificación de correctitud}
La opción \texttt{-{}-verify} de \texttt{lu\_distributed.py} implementa la verificación. Con $N=1024$, $b=128$ y los \emph{grids} $1\times1$, $1\times2$, $2\times1$ y $2\times2$ (\texttt{run\_experimento\_b.py}), al terminar la factorización se juntan los arreglos locales en el rank 0 (\texttt{comm.gather}), se reconstruye la matriz global deshaciendo la distribución bloque-cíclica y se compara contra:
\begin{enumerate}
  \item el resultado de \texttt{lu\_serial.py} sobre la misma matriz, mediante la diferencia relativa \[\max_{ij}|A^{\text{dist}}_{ij}-A^{\text{serial}}_{ij}|\,/\,\max_{ij}|A^{\text{serial}}_{ij}|;\]
  \item la definición, mediante el residuo relativo $\|A-LU\|_F/\|A\|_F$, con $L$ (triángulo inferior con unos en la diagonal) y $U$ extraídos de la matriz factorizada.
\end{enumerate}
También se cuenta en cuántos bloques diagonales \texttt{lu\_factor} permutó filas (\texttt{piv} distinto de $0,1,\dots,b-1$), ya que la implementación descarta el vector de pivoteo; este número debe ser 0. La versión serial se valida aparte con el mismo residuo.

\begin{table}[H]
\centering\small
\caption{Verificación de correctitud con $N=1024$, $b=128$ en el computador personal.}
\label{tab:verif}
\begin{tabular}{lccc}
\toprule
\textbf{Grid} & Dif. relativa máx. vs.\ serial & $\|A-LU\|_F/\|A\|_F$ & Bloques con pivoteo \\
\midrule
Serial & -- & $4.2\times10^{-16}$ & 0 \\
$1\times1$ & 0 & $4.2\times10^{-16}$ & 0 \\
$1\times2$ & 0 & $4.2\times10^{-16}$ & 0 \\
$2\times1$ & 0 & $4.2\times10^{-16}$ & 0 \\
$2\times2$ & 0 & $4.2\times10^{-16}$ & 0 \\
\bottomrule
\end{tabular}
\end{table}

En los cuatro \emph{grids} la factorización distribuida es idéntica a la serial (diferencia exactamente 0): cada elemento se obtiene con la misma secuencia de operaciones, solo que repartida entre procesos. El residuo es $4.2\times10^{-16}$, del orden de la precisión de máquina ($\varepsilon\approx1.1\times10^{-16}$), y \texttt{lu\_factor} nunca permutó filas, lo que confirma que descartar el vector de pivoteo es válido para esta matriz.

\subsection{Prueba local}
Con $N=2048$, $b=128$ y \emph{grid} $2\times2$ en el computador personal (4 procesos en 10 cores, sin sobresuscripción) se obtuvieron los tiempos de cada proceso de la Tabla~\ref{tab:blocal}.

\begin{table}[H]
\centering\small
\caption{Tiempos por proceso con $N=2048$, $b=128$, \emph{grid} $2\times2$ y $T=1$ en el computador personal (mediana de 3).}
\label{tab:blocal}
\begin{tabular}{ccrrr}
\toprule
Rank & Posición & $\Tcomp$ [ms] & $\Tcomm$ [ms] & $\Ttot$ [ms] \\
\midrule
0 & $(0,0)$ & 17.6 & 6.9 & 24.7 \\
1 & $(0,1)$ & 18.9 & 5.9 & 24.8 \\
2 & $(1,0)$ & 18.9 & 5.9 & 24.8 \\
3 & $(1,1)$ & 19.7 & 5.1 & 24.8 \\
\bottomrule
\end{tabular}
\end{table}

El cómputo está bien balanceado (17.6 a 19.7 ms, un 11\% de diferencia) y la comunicación representa entre el 21\% y el 28\% del tiempo. La suma $\Tcomp+\Tcomm$ es casi igual en los cuatro procesos: el que calcula menos espera más, porque esa espera queda dentro de los \texttt{Bcast}. El rank 0 es el que menos calcula: tiene las filas y columnas de bloque pares ($0,2,\dots,14$), así que en cada paso le quedan igual o menos bloques en la submatriz restante que a los demás, y es el que más espera. El rank 3 tiene las impares ($1,3,\dots,15$) y es el caso opuesto.

Esta prueba sirve para validar la medición, no para estudiar el escalamiento: el problema es pequeño (5.7 GFLOP en total) y los 4 procesos juntos alcanzan unos 230 GFLOP/s, similar a un solo \emph{thread} en la parte (a). Una explicación probable es que en los chips de Apple la unidad de matrices es compartida por los cores de un mismo clúster, de modo que varios procesos compiten por ella. Por eso el escalamiento se estudia en el \emph{cluster}.

% ======================================================================
\section{(c) Tamaño de bloque en el \emph{cluster}}
Configuración 2 ($2\times2$, $T=2$, un nodo, 4 procesos), $N=2^{13}=8192$.

\begin{figure}[H]\centering
\begin{tikzpicture}
\begin{axis}[
  width=0.8\linewidth, height=6.5cm,
  ybar stacked, bar width=14pt,
  symbolic x coords={64,128,256,512,1024,2048,4096},
  xtick=data, xlabel={Tamaño de bloque $b$}, ylabel={Tiempo rank 0 [s]},
  title={$N=8192$, grid $2\times2$, $T=2$ (DATOS PROVISORIOS)},
  legend style={at={(0.5,-0.22)}, anchor=north, legend columns=-1, font=\small},
  ymin=0, grid=major, label style={font=\small}, tick label style={font=\small}]
\addplot[fill=blue!50] coordinates {(64,6.5379) (128,3.2987) (256,2.6469) (512,2.4230) (1024,2.4226) (2048,2.8455) (4096,4.7042)};
\addplot[fill=orange!70] coordinates {(64,5.1950) (128,2.3568) (256,1.4850) (512,1.0186) (1024,0.8184) (2048,1.5115) (4096,1.2054)};
\legend{$\Tcomp$, $\Tcomm$}
\end{axis}
\end{tikzpicture}
\caption{Tiempo de cómputo y de comunicación del rank 0 según $b$. \falta{Reemplazar con los datos corregidos.}}
\label{fig:c}
\end{figure}

\begin{table}[H]
\centering\small
\caption{Barrido de $b$ con $N=8192$ en la configuración 2 ($2\times2$, $T=2$). \falta{Datos provisorios: ver nota.}}
\label{tab:c}
\begin{tabular}{rrrrrr}
\toprule
$b$ & $N/b$ & $\Tcomp$ [s] & $\Tcomm$ [s] & $\Ttot$ [s] & GFLOP/s \\
\midrule
8192 & 1   & \multicolumn{4}{c}{sin memoria (\emph{OOM}, límite 4 GB)} \\
4096 & 2   & 4.70$^\dagger$ & 1.21$^\dagger$ & 16.39$^\dagger$ & 22.4$^\dagger$ \\
2048 & 4   & 2.85 & 1.51 & 5.98 & 61.3 \\
1024 & 8   & 2.42 & 0.82 & 3.55 & 103.2 \\
\textbf{512} & \textbf{16} & \textbf{2.42} & \textbf{1.02} & \textbf{3.53} & \textbf{103.8} \\
256  & 32  & 2.65 & 1.49 & 4.16 & 88.1 \\
128  & 64  & 3.30 & 2.36 & 5.67 & 64.6 \\
64   & 128 & 6.54 & 5.20 & 11.75 & 31.2 \\
\bottomrule
\end{tabular}\\[2pt]
{\footnotesize $^\dagger$ Corrida en 2 nodos con $T=1$ (otra configuración) porque en 1 nodo fue terminada por memoria.}
\end{table}

\nota{\textbf{Estos datos no son válidos todavía}: el código distribuido no aplica la actualización en los procesos que no son dueños de $A_{kk}$, así que hace mucho menos trabajo que el real y la matriz resultante es incorrecta. Con el código corregido hay que repetir todo el barrido. Lo que sigue es la estructura del análisis con la tendencia actual, a confirmar.}

Con los datos provisorios, el mínimo está en $b=512$ (3.53 s), prácticamente empatado con $b=1024$ (3.55 s), que coincide con $b_a=512$ de la parte (a). \falta{confirmar con los datos nuevos.} Las razones de que el óptimo se desplace hacia bloques más grandes son:
\begin{itemize}
  \item \textbf{Comunicación:} en la versión distribuida aparece un costo por mensaje (latencia $\alpha$) que no existe en la serial. El número de \texttt{Bcast} crece como $(N/b)^2$, y con $b=64$ la comunicación ya es casi igual al cómputo (5.20 frente a 6.54 s).
  \item \textbf{Eficiencia de la actualización:} con $b$ pequeño la actualización de rango $b$ tiene baja intensidad aritmética y queda limitada por memoria, igual que en (a); además hay $N/b$ pasos, cada uno con varias llamadas desde Python.
  \item \textbf{Hardware:} el \emph{cluster} tiene núcleos con cachés más grandes y, con $T=2$, BLAS reparte cada producto entre 2 threads, lo que favorece bloques más grandes. \falta{completar con cachés del cluster}.
  \item \textbf{Bloques grandes:} con $b\ge2048$ hay $N/b\le4$ bloques por dimensión y solo uno o dos bloques por proceso; la mayoría de los procesos queda inactiva en cada paso (ver parte f), y con $b=N$ todo el trabajo y toda la memoria recaen en un solo proceso, que excede los 4 GB del nodo.
\end{itemize}
Por lo tanto, para las partes siguientes se usan $b_a=512$ (parte a) y $b_c=\falta{512}$ (parte c). Si ambos coinciden, se agrega $b=128$ como contraste, para poder comparar dos tamaños de bloque en las partes (d) y (e).

% ======================================================================
\section{(d) Escalabilidad en un nodo: configuraciones 0 a 3}
\textbf{Elección de $T(1)$.} Como tiempo base se usa la configuración 0 con $T=1$ ($1\times1$, un proceso, un thread), ejecutando el \emph{mismo} programa \texttt{lu\_distributed.py} con el mismo $b$. Así $T(1)$ mide exactamente el mismo algoritmo y la misma implementación en un solo CPU, sin comunicación (todos los \texttt{Bcast} son sobre comunicadores de tamaño 1), y el speedup refleja solo el efecto de agregar CPUs. Se calcula $S=T(1)/T$ y $E=S/(p\cdot T)$, donde $p\cdot T$ es el total de CPUs usados.
\nota{Alternativa a discutir: usar el mejor de los dos $T(1)$ (con $b_a$ o con $b_c$) como base común para ambas columnas, para que el speedup sea comparable entre valores de $b$. Decidir y justificar.}

\begin{table}[H]
\centering\small
\caption{Configuraciones 0--3 en un nodo, $N=8192$. \falta{Completar con las corridas corregidas.}}
\label{tab:d}
\begin{tabular}{cccc|rrr|rrr}
\toprule
 & & & & \multicolumn{3}{c|}{$b=b_a=512$} & \multicolumn{3}{c}{$b=\falta{128 o $b_c$}$} \\
Conf. & Grid & $T$ & CPUs & $\Ttot$ [s] & $S$ & $E$ & $\Ttot$ [s] & $S$ & $E$ \\
\midrule
0 & $1\times1$ & 1 & 1 & \falta{} & 1.00 & 1.00 & \falta{} & 1.00 & 1.00 \\
0 & $1\times1$ & 2 & 2 & \falta{} & & & \falta{} & & \\
0 & $1\times1$ & 4 & 4 & \falta{} & & & \falta{} & & \\
0 & $1\times1$ & 8 & 8 & \falta{} & & & \falta{} & & \\
\midrule
1 & $1\times2$ & 1 & 2 & \falta{} & & & \falta{} & & \\
1 & $1\times2$ & 2 & 4 & \falta{} & & & \falta{} & & \\
1 & $1\times2$ & 4 & 8 & \falta{} & & & \falta{} & & \\
1 & $2\times1$ & 1 & 2 & \falta{} & & & \falta{} & & \\
1 & $2\times1$ & 2 & 4 & \falta{} & & & \falta{} & & \\
1 & $2\times1$ & 4 & 8 & \falta{} & & & \falta{} & & \\
\midrule
2 & $2\times2$ & 1 & 4 & \falta{} & & & \falta{} & & \\
2 & $2\times2$ & 2 & 8 & \falta{} & & & \falta{} & & \\
\midrule
3 & $2\times4$ & 1 & 8 & \falta{} & & & \falta{} & & \\
3 & $4\times2$ & 1 & 8 & \falta{} & & & \falta{} & & \\
\bottomrule
\end{tabular}
\end{table}

\falta{Figura: speedup frente a CPUs totales (1, 2, 4, 8), una curva por configuración y la línea ideal $S=p\cdot T$; dos paneles, uno por cada $b$.}

\textbf{Análisis (a completar con los datos).} Preguntas a responder con la tabla:
\begin{itemize}
  \item \textbf{Mejor configuración con 8 CPUs:} comparar $1\times1$ con $T=8$ (todo el paralelismo en BLAS, sin comunicación), $1\times2$/$2\times1$ con $T=4$, $2\times2$ con $T=2$ y $2\times4$/$4\times2$ con $T=1$. Se espera que más procesos MPI con menos threads gane cuando $b$ es pequeño (los productos de $b\times b$ son demasiado pequeños para que BLAS aproveche varios threads), y que más threads gane cuando $b$ es grande. \falta{conclusión}.
  \item \textbf{Forma del grid:} con $p$ fijo, el volumen de comunicación de los paneles por proceso es $\sim \tfrac{N^2}{2}\big(\tfrac{\log_2 p_r}{p_c}+\tfrac{\log_2 p_c}{p_r}\big)$ palabras (parte g), que es simétrico en $p_r$ y $p_c$. La asimetría viene de la ruta crítica: los $p_c$ procesos de la fila $r$ calculan el panel $U$ y los $p_r$ de la columna $c$ calculan el panel $L$, y el trabajo del panel por proceso es $\sim\frac{m}{p_c}+\frac{m}{p_r}$ bloques. \falta{comparar $1\times2$ vs $2\times1$ y $2\times4$ vs $4\times2$ y concluir si la forma afecta}.
\end{itemize}
\nota{Los datos de \texttt{results\_d.md} no sirven para esta parte: $1\times2$ y $2\times2$ se corrieron repartidos en 2 nodos (\texttt{-{}-map-by ppr:1:node} / \texttt{ppr:2:node}), $2\times4$ también en 2 nodos (eso es la configuración 4, no la 3), no se varió $T$ y no se corrió $2\times1$ ni $4\times2$ en un nodo. Además el bug de la actualización hace que los speedups salgan inflados.}

% ======================================================================
\section{(e) Uno frente a dos nodos: configuraciones 3, 4 y 5}
Con $N=8192$ se compara el \emph{grid} $2\times4$ en un nodo con 8 procesos (configuración 3) y en dos nodos con 4 procesos cada uno (configuración 4). Ambas usan exactamente los mismos procesos, los mismos bloques y los mismos mensajes; solo cambia el medio por donde viajan los mensajes que cruzan de un nodo a otro. La configuración 5 ($4\times4$, dos nodos, 16 CPUs) muestra qué pasa al duplicar los procesos.

\begin{table}[H]
\centering\small
\caption{Configuraciones 3, 4 y 5 con $N=8192$, $T=1$. Tiempos máximos sobre los procesos. \falta{Completar.}}
\label{tab:e}
\begin{tabular}{ccccc|rrr|rrr}
\toprule
 & & & & & \multicolumn{3}{c|}{$b=128$} & \multicolumn{3}{c}{$b=512$} \\
Conf. & Grid & Nodos & Proc./nodo & CPUs & $\Tcomp$ & $\Tcomm$ & $\Ttot$ & $\Tcomp$ & $\Tcomm$ & $\Ttot$ \\
\midrule
3 & $2\times4$ & 1 & 8 & 8  & \falta{} & \falta{} & \falta{} & \falta{} & \falta{} & \falta{} \\
4 & $2\times4$ & 2 & 4 & 8  & \falta{} & \falta{} & \falta{} & \falta{} & \falta{} & \falta{} \\
5 & $4\times4$ & 2 & 8 & 16 & \falta{} & \falta{} & \falta{} & \falta{} & \falta{} & \falta{} \\
\bottomrule
\end{tabular}
\end{table}

\falta{Figura: barras apiladas $\Tcomp$/$\Tcomm$ para las configuraciones 3, 4 y 5, un grupo por cada $b$.}

\textbf{Análisis.} El cómputo debería ser casi igual en las configuraciones 3 y 4, porque cada proceso tiene los mismos bloques y hace las mismas operaciones; la diferencia debe aparecer en $\Tcomm$. Dentro de un nodo, Open MPI usa memoria compartida: la latencia es del orden de $\sim$1 $\mu$s y el ancho de banda está limitado por la copia en memoria (decenas de GB/s). Entre nodos, cada mensaje pasa por la red, con una latencia $\alpha$ de \falta{decenas de $\mu$s en Ethernet / $\sim$1--2 $\mu$s en InfiniBand} y un ancho de banda $1/\beta$ de \falta{x GB/s}. Con el \emph{grid} $2\times4$ repartido en 2 nodos de 4 procesos (ranks 0--3 en un nodo, ranks 4--7 en otro), cada fila del \emph{grid} queda dentro de un nodo, pero cada columna cruza la red: los \texttt{Bcast} de $U_{kj}$ en \texttt{col\_comm} son inter-nodo y los de $L_{ik}$ en \texttt{row\_comm} son intra-nodo. \falta{verificar el mapeo de ranks a nodos con \texttt{-{}-report-bindings} o imprimiendo \texttt{MPI.Get\_processor\_name()}}.

La dependencia con $b$ separa los dos efectos. Con un \texttt{Bcast} por panel, hay $\sim2N/b$ difusiones en la ruta crítica, mientras que el volumen total de los paneles ($\sim\tfrac{N^2}{2}(\tfrac{\log_2p_r}{p_c}+\tfrac{\log_2p_c}{p_r})$ palabras) casi no depende de $b$ (parte g). Con $b=128$ hay 128 pasos con mensajes más pequeños, donde pesa más la \textbf{latencia}; con $b=512$ hay solo 32 pasos con mensajes 4 veces más grandes, donde pesa más el \textbf{ancho de banda}. \falta{Cuantificar: estimar $\alpha$ y $\beta$ inter-nodo a partir de la diferencia $\Tcomm^{(4)}-\Tcomm^{(3)}$ para ambos $b$.} En la configuración 5, $\Tcomp$ debería caer a la mitad, pero $\Tcomm$ crece: cada \texttt{Bcast} por columna tiene más pasos ($\log_2 4=2$ en vez de $\log_2 2=1$) y, con las filas 0--1 en un nodo y 2--3 en el otro, todos los \texttt{col\_comm} cruzan la red, ahora entre 4 procesos.

\nota{Dato provisorio con el código con bug: la configuración 4 ($2\times4$, dos nodos, $b=512$) dio $\Ttot = 2.97$, $3.60$ y $3.48$ s en tres corridas, y la 5 dio $\Ttot=3.44$ s con $\Tcomm=2.56$ s (74\% del total). Nunca se corrió la configuración 3, que es la comparación central de esta pregunta. La comparación de formas $1\times8$/$2\times4$/$4\times2$/$8\times1$ de \texttt{results\_e.md} no la pide el enunciado; se puede dejar en un anexo como experimento adicional.}

% ======================================================================
\section{(f) Bloques grandes frente a bloques pequeños en el \emph{grid} $2\times2$}
Se compara $b=N/2=4096$ con $b=N/16=512$ en el \emph{grid} $2\times2$, con $N=8192$, en un nodo con $T=2$ (configuración 2), de modo que ambos casos salen del mismo barrido de la parte (c).

\begin{table}[H]
\centering\small
\caption{Tiempos de cada proceso en el \emph{grid} $2\times2$, $N=8192$. \falta{Completar.}}
\label{tab:f}
\begin{tabular}{cc|rrr|rrr}
\toprule
 & & \multicolumn{3}{c|}{$b=N/2=4096$} & \multicolumn{3}{c}{$b=N/16=512$} \\
Rank & Posición & $\Tcomp$ [s] & $\Tcomm$ [s] & $\Ttot$ [s] & $\Tcomp$ [s] & $\Tcomm$ [s] & $\Ttot$ [s] \\
\midrule
0 & $(0,0)$ & \falta{} & \falta{} & \falta{} & \falta{} & \falta{} & \falta{} \\
1 & $(0,1)$ & \falta{} & \falta{} & \falta{} & \falta{} & \falta{} & \falta{} \\
2 & $(1,0)$ & \falta{} & \falta{} & \falta{} & \falta{} & \falta{} & \falta{} \\
3 & $(1,1)$ & \falta{} & \falta{} & \falta{} & \falta{} & \falta{} & \falta{} \\
\bottomrule
\end{tabular}
\end{table}

\falta{Figura: barras apiladas $\Tcomp$/$\Tcomm$ por rank, dos paneles ($b=4096$ y $b=512$).}

\textbf{Qué se espera y por qué.} Con $b=N/2$ la matriz tiene $2\times2$ bloques y cada proceso guarda \emph{exactamente uno}. El algoritmo queda completamente secuencial:
\begin{enumerate}
  \item Paso $k=0$: $(0,0)$ factoriza $A_{00}$ ($\tfrac{2}{3}b^3$ FLOP) mientras los demás esperan; luego $(0,1)$ calcula $U_{01}$ y $(1,0)$ calcula $L_{10}$ ($b^3$ cada uno, en paralelo); finalmente solo $(1,1)$ actualiza $A_{11}\leftarrow A_{11}-L_{10}U_{01}$ ($2b^3$).
  \item Paso $k=1$: solo $(1,1)$ factoriza $A_{11}$ ($\tfrac{2}{3}b^3$).
\end{enumerate}
La ruta crítica suma $\tfrac{2}{3}+1+2+\tfrac{2}{3}=\tfrac{13}{3}b^3$ FLOP, frente a $\tfrac{16}{3}b^3$ del total: el máximo speedup posible es $16/13\approx1.23$ con 4 procesos. Además, cada mensaje es un bloque de $4096^2\times8$ B $=128$ MB, y en el paso 0 se difunden tres bloques así ($A_{00}$, $U_{01}$ y $L_{10}$). Se espera entonces que $(1,1)$ tenga la mayor parte de $\Tcomp$, y que los demás procesos acumulen tiempo de espera, que aparece en $\Tcomm$ (esperando a la raíz en el \texttt{Bcast}) o como diferencia entre $\Ttot$ y $\Tcomp+\Tcomm$ (esperando en el \texttt{Barrier} final).

Con $b=N/16$ hay $16\times16$ bloques, $8\times8$ por proceso, y la distribución cíclica reparte la submatriz restante de forma pareja en casi todos los pasos: el trabajo se balancea y los mensajes son de 2 MB. \falta{Contrastar con la tabla: desbalance de $\Tcomp$ entre ranks, peso de $\Tcomm$.}

\nota{Dato provisorio (con bug, solo rank 0): con $b=4096$ en 2 nodos y $T=1$, $\Tcomp=4.70$, $\Tcomm=1.21$ y $\Ttot=16.39$ s; con $b=512$ en 1 nodo y $T=2$, $3.53$ s. El hueco $\Ttot-\Tcomp-\Tcomm\approx10$ s en $b=4096$ es justamente el rank 0 esperando en el \texttt{Barrier} final a los otros procesos, que coincide con la explicación de arriba. Pero ambas corridas usan configuraciones distintas y no son comparables: hay que repetir ambas en la misma configuración. La tabla de \texttt{results\_f.md} ($2\times4$ con $b\in\{2048,512,128\}$) no responde esta pregunta.}

% ======================================================================
\section{(g) Modelo de costo}
Se modela cada mensaje de $n$ palabras como $\alpha+n\beta$, cada FLOP como $\gamma$, y un \texttt{Bcast} entre $q$ procesos con un árbol binomial como $\lceil\log_2 q\rceil(\alpha+n\beta)$. Sea $n_b=N/b$ y, en el paso $k$, $m=n_b-k-1$ el número de filas (y columnas) de bloques restantes. Con distribución bloque-cíclica, cada proceso tiene a lo más $\lceil m/p_r\rceil$ filas y $\lceil m/p_c\rceil$ columnas de bloques restantes. Siguiendo la implementación (un \texttt{Bcast} por panel), el costo de la ruta crítica en el paso $k$ es:

\begin{center}
\small
\begin{tabular}{lll}
\toprule
\textbf{Etapa} & \textbf{Cómputo} & \textbf{Comunicación} \\
\midrule
Factorizar $A_{kk}$ & $\tfrac{2}{3}b^3\gamma$ & -- \\
Difundir $A_{kk}$ & -- & $(\log_2 p_c+\log_2 p_r)(\alpha+b^2\beta)$ \\
Paneles $U$ y $L$ & $\big(\tfrac{m}{p_c}+\tfrac{m}{p_r}\big)b^3\gamma$ & -- \\
Difundir paneles & -- & $\log_2 p_r\big(\alpha+\tfrac{m}{p_c}b^2\beta\big)+\log_2 p_c\big(\alpha+\tfrac{m}{p_r}b^2\beta\big)$ \\
Actualización & $\tfrac{m}{p_r}\tfrac{m}{p_c}\,2b^3\gamma$ & -- \\
\bottomrule
\end{tabular}
\end{center}

Sumando sobre los $n_b$ pasos, con $\sum m\approx n_b^2/2$ y $\sum m^2\approx n_b^3/3$:
\begin{align*}
T(N,b,p_r,p_c) \;\approx\;&
\underbrace{\frac{2N^3}{3p}\gamma}_{\text{actualización}}
+\underbrace{\frac{N^2b}{2}\Big(\frac{1}{p_r}+\frac{1}{p_c}\Big)\gamma}_{\text{paneles}}
+\underbrace{\frac{2}{3}Nb^2\gamma}_{\text{bloques diagonales}}\\[2pt]
&+\underbrace{\frac{2N}{b}\log_2 p\;\alpha}_{\text{latencia}}
+\underbrace{\Big[Nb\log_2 p+\frac{N^2}{2}\Big(\frac{\log_2 p_r}{p_c}+\frac{\log_2 p_c}{p_r}\Big)\Big]\beta}_{\text{ancho de banda}}
\end{align*}
con $p=p_rp_c$.

\textbf{Qué explica el $b$ óptimo.} El término principal $\tfrac{2N^3}{3p}\gamma$ no depende de $b$, y el volumen de datos de los paneles ($\beta N^2/2\cdots$) tampoco. Los términos que \emph{crecen} con $b$ son los de la ruta crítica secuencial: los paneles ($\propto N^2b$), que en cada paso calcula solo una fila y una columna del \emph{grid}, y los bloques diagonales ($\propto Nb^2$), que calcula un solo proceso mientras los demás esperan. Ambos representan también el desbalance: cuando $m<p_r$ o $m<p_c$ hay procesos sin trabajo, y con $b=N/2$ el algoritmo es casi secuencial (parte f). El término que \emph{decrece} con $b$ es la latencia, proporcional al número de pasos $N/b$. En Python hay que sumarle el costo fijo $\delta$ de las $\sim5$ llamadas a NumPy/SciPy/mpi4py de cada paso, que tiene la misma forma: $\tfrac{N}{b}(2\alpha\log_2 p+5\delta)$. Igualando las derivadas del término de paneles y de este término:
\[
b_{\text{opt}}\approx\sqrt{\frac{2\,(2\alpha\log_2 p+5\delta)}{N\gamma\,(1/p_r+1/p_c)}}.
\]
Con valores típicos ($\alpha,\delta\sim10^{-5}$ s, $\gamma\sim(25\ \text{GFLOP/s})^{-1}$), esta expresión da $b_{\text{opt}}$ de unas pocas decenas, mucho menor que lo medido. La diferencia la explica un efecto que el modelo con $\gamma$ constante no captura: \textbf{$\gamma$ depende de $b$}. La actualización es una multiplicación de rango $b$, $(m b\times b)\cdot(b\times mb)$, cuya intensidad aritmética es $\sim b/8$ FLOP por byte (cada elemento de la submatriz, 16 bytes leídos y escritos, recibe $2b$ FLOP): con $b$ pequeño BLAS queda limitado por el ancho de banda de memoria y $\gamma$ efectivo sube mucho. Lo mismo ocurre con \texttt{lu\_factor} y \texttt{solve\_triangular} sobre bloques pequeños. Por eso el óptimo real queda donde la multiplicación ya alcanza su rendimiento máximo ($b$ de algunos cientos) y antes de que los términos de ruta crítica y desbalance crezcan.
\falta{Con las mediciones: estimar $\gamma(b)$ de la configuración 0 con $T=1$ ($\gamma\approx\Ttot/(\tfrac{2}{3}N^3)$), $\alpha$ y $\beta$ inter-nodo de la parte (e), evaluar el modelo y compararlo con el $b_{\text{opt}}$ medido en (c).}

\textbf{Cuello de botella.} \falta{Responder con los datos corregidos.} Con los datos provisorios, $\Tcomm$ es del mismo orden que $\Tcomp$ desde 4 procesos y lo supera en dos nodos (por ejemplo, 74\% del tiempo en $4\times4$), lo que apunta a la comunicación (y la espera que incluye) como cuello de botella al escalar. Como el bug reduce artificialmente $\Tcomp$, esta conclusión puede cambiar.

% ======================================================================
\section{Conclusiones}
\nota{Escribir con los datos corregidos. Ideas: (1) el $b$ óptimo es un compromiso entre eficiencia de BLAS / overhead de Python (favorece $b$ grande) y latencia + ruta crítica + balance de carga (favorece $b$ pequeño), y se desplaza hacia arriba al pasar a MPI; (2) qué combinación procesos $\times$ threads gana con 8 CPUs; (3) si la forma del grid importa; (4) cuánto cuesta cruzar la red entre nodos; (5) $b=N/2$ elimina el paralelismo; (6) cuello de botella final.}

% ======================================================================
\section*{Declaración de uso de IA}
Para este trabajo se usó Claude (Anthropic) mediante Claude Code como asistente para redactar la estructura del informe y el desarrollo del modelo de costo, y para revisar el código y los resultados entregados por el grupo, a partir de las mediciones realizadas por los integrantes en su computador y en el \emph{cluster}. Los resultados corresponden a mediciones reales y el análisis final fue revisado por los integrantes.
\nota{Ajustar según cómo usaron IA para escribir y correr el código.}

\appendix
\section{Scripts SLURM}
Los scripts se entregan en el \texttt{.zip}. Todos incluyen \texttt{-{}-qos}, \texttt{-{}-partition} y \texttt{-{}-account} de la partición del curso, fijan las variables de \emph{threads} en \texttt{\$SLURM\_CPUS\_PER\_TASK} y usan \texttt{-{}-ntasks-per-node} para controlar la ubicación de los procesos.
\begin{table}[H]
\centering\small
\caption{Scripts SLURM (carpeta \texttt{slurm/}). Todos usan $N=8192$ y 3 repeticiones; \texttt{comun.sh} define la función que lanza cada corrida.}
\begin{tabular}{llp{9cm}}
\toprule
\textbf{Script} & \textbf{Partes} & \textbf{Corridas} \\
\midrule
\texttt{job1\_parte\_c\_f} & c, f & 1 nodo, 4 tareas $\times$ 2 CPUs: \emph{grid} $2\times2$, $T=2$, $b=4096,\dots,64$ (y $b=N$). \\
\texttt{job2\_parte\_d\_T1} & d, e & 1 nodo, 8 tareas $\times$ 1 CPU: $1\times1$, $1\times2$, $2\times1$, $2\times2$, $2\times4$ y $4\times2$ con $T=1$, para $b_a$ y $b_c$. \\
\texttt{job3\_parte\_d\_Tvar} & d & 1 nodo, 2 tareas $\times$ 4 CPUs: $1\times1$ con $T=2,4,8$; $1\times2$ y $2\times1$ con $T=2,4$, para $b_a$ y $b_c$. Cada corrida usa $p\cdot T\le8$ CPUs. \\
\texttt{job4\_parte\_e} & e & 2 nodos, 8 tareas $\times$ 1 CPU por nodo: $2\times4$ con 4 procesos por nodo y $4\times4$ con 8 por nodo, $T=1$, para $b_a$ y $b_c$. \\
\bottomrule
\end{tabular}
\end{table}

\end{document}
